\documentclass[12pt]{article}
\usepackage[UTF8]{inputenc}
\usepackage{thaisupp}
\usepackage{geometry}
\geometry{a4paper, margin=2.5cm}
\begin{document}
\begin{center}
{\LARGE\textbf{ฟิสิกส์ ม.ปลาย — Part 1}}\\[0.3em]
{\Large\textbf{Lesson 10: เสียง (Extra) (Extra)}}\\[0.5em]
\rule{0.8\linewidth}{0.5pt}
\end{center}
\section*{คำถาม In-Class}
\begin{enumerate}
\item \textbf{อัตราเร็วเสียงในอากาศที่ 0°C เป็นเท่าใด}
\begin{enumerate}
\item ก) 300 m/s
\item ข) 331 m/s
\item ค) 340 m/s
\item ง) 350 m/s
\end{enumerate}
\item \textbf{ปรากฏการณ์ดอปเปลอร์เกี่ยวข้องกับการเปลี่ยนแปลงของข้อใด}
\begin{enumerate}
\item ก) ความยาวคลื่น
\item ข) ความถี่
\item ค) อัตราเร็ว
\item ง) ทิศทาง
\end{enumerate}
\item \textbf{เสียงที่มีระดับเสียง 0 dB มีความเข้มเท่ากับเท่าใด}
\begin{enumerate}
\item ก) 0 W/m²
\item ข) 10⁻¹² W/m²
\item ค) 10⁻⁶ W/m²
\item ง) 1 W/m²
\end{enumerate}
\item \textbf{แหล่งเสียงความถี่ 400 Hz เคลื่อนที่ด้วยความเร็ว 50 m/s เข้าหาผู้สังเกตที่อยู่กับที่ ความถี่ที่ได้ยินเป็นเท่าใด (v = 340 m/s)}
\begin{enumerate}
\item ก) 480 Hz
\item ข) 545 Hz
\item ค) 500 Hz
\item ง) 476 Hz
\end{enumerate}
\item \textbf{เสียงต่ำ (low pitch) หมายถึงเสียงที่มีความถี่อย่างไร}
\begin{enumerate}
\item ก) ต่ำ
\item ข) สูง
\item ค) ปานกลาง
\item ง) คงที่
\end{enumerate}
\item \textbf{ระดับเสียงเพิ่มขึ้นจาก 30 dB เป็น 50 dB ความเข้มเพิ่มขึ้นกี่เท่า}
\begin{enumerate}
\item ก) 20 เท่า
\item ข) 100 เท่า
\item ค) 10 เท่า
\item ง) 200 เท่า
\end{enumerate}
\item \textbf{รถไฟเสียงนกหวีดความถี่ 440 Hz วิ่งเข้าหาอุโมงค์ด้วยความเร็ว 30 m/s ความถี่ที่ผู้โดยสารในรถไฟได้ยินเป็นเท่าใด (v = 340 m/s)}
\begin{enumerate}
\item ก) 440 Hz
\item ข) 380 Hz
\item ค) 478 Hz
\item ง) 400 Hz
\end{enumerate}
\item \textbf{คลื่นเสียงในอากาศเป็นคลื่นชนิดใด}
\begin{enumerate}
\item ก) คลื่นตามขวาง
\item ข) คลื่นตามยาว
\item ค) คลื่นผิวน้ำ
\item ง) คลื่นแม่เหล็กไฟฟ้า
\end{enumerate}
\item \textbf{ถ้าความเข้มเสียงเพิ่มขึ้น 10,000 เท่า ระดับเสียงเพิ่มขึ้นกี่ dB}
\begin{enumerate}
\item ก) 20 dB
\item ข) 30 dB
\item ค) 40 dB
\item ง) 50 dB
\end{enumerate}
\item \textbf{รถพยุมยางเสียงดัง 600 Hz เคลื่อนที่เข้าหาผู้สังเกตด้วยความเร็ว 25 m/s และผู้สังเกตเคลื่อนที่เข้าหารถด้วยความเร็ว 5 m/s ความถี่ที่ได้ยินเป็นเท่าใด (v = 340 m/s)}
\begin{enumerate}
\item ก) 620 Hz
\item ข) 655 Hz
\item ค) 700 Hz
\item ง) 720 Hz
\end{enumerate}
\item \textbf{เสียงอินฟราซาวด์มีความถี่ต่ำกว่าเท่าใด}
\begin{enumerate}
\item ก) 2 Hz
\item ข) 20 Hz
\item ค) 200 Hz
\item ง) 2000 Hz
\end{enumerate}
\item \textbf{เสียงที่มีระดับดัง 140 dB สามารถทำให้เกิดอาการบาดเจ็บอะไรได้}
\begin{enumerate}
\item ก) ปวดหู
\item ข) หูหนวกชั่วคราว
\item ค) แก้วหูฉีกขาด
\item ง) ทั้งหมดที่กล่าวมา
\end{enumerate}
\item \textbf{เครื่องบินเจ็ตบินด้วยความเร็วเสียง (Mach 1) ทำให้เกิดเสียงอะไร}
\begin{enumerate}
\item ก) เสียงปกติ
\item ข) เสียงดังมาก
\item ค) เสียงโซนิกบูม
\item ง) ความถี่เพิ่มเป็น 2 เท่า
\end{enumerate}
\item \textbf{อัตราเร็วเสียงในน้ำประมาณเท่าใด}
\begin{enumerate}
\item ก) 340 m/s
\item ข) 750 m/s
\item ค) 1500 m/s
\item ง) 3000 m/s
\end{enumerate}
\item \textbf{เมื่อรถพยุมยางเสียงดังวิ่งผ่านผู้สังเกตที่อยู่กับที่ เสียงที่ได้ยินจะเปลี่ยนแปลงอย่างไร}
\begin{enumerate}
\item ก) ดังขึ้นตลอด
\item ข) เบาลงตลอด
\item ค) แหลมสูงขึ้นแล้วทันใดก็ต่ำลง
\item ง) ความถี่คงที่ตลอด
\end{enumerate}
\item \textbf{ความถี่ 440 Hz เป็นเสียงที่ใช้เป็นมาตรฐานการลงครึ่งเสียง เรียกว่าอะไร}
\begin{enumerate}
\item ก) เสียง A กลาง
\item ข) เสียง A มาตรฐาน
\item ค) เสียง C กลาง
\item ง) เสียง D มาตรฐาน
\end{enumerate}
\item \textbf{เสียงเดินทางในอากาศได้เร็วขึ้นเมื่ออุณหภูมิสูงขึ้น เพราะเหตุใด}
\begin{enumerate}
\item ก) อากาศร้อนขยายตัว
\item ข) โมเลกุลอากาศเคลื่อนที่เร็วขึ้น ส่งผ่านการสั่นได้เร็วขึ้น
\item ค) อากาศร้อนมีความหนาแน่นมากขึ้น
\item ง) อากาศร้อนมีความชื้นมากขึ้น
\end{enumerate}
\item \textbf{ระดับเสียงของเครื่องยนต์จรวด 180 dB มีความเข้มเป็นกี่เท่าของเสียง 0 dB}
\begin{enumerate}
\item ก) 10¹⁸ เท่า
\item ข) 10⁸ เท่า
\item ค) 10⁶ เท่า
\item ง) 10⁴ เท่า
\end{enumerate}
\item \textbf{สูตรปรากฏการณ์ดอปเปลอร์ f' = f(v/(v-vs)) ใช้เมื่อใด}
\begin{enumerate}
\item ก) เมื่อแหล่งเสียงอยู่กับที่ ผู้สังเกตเคลื่อนที่
\item ข) เมื่อแหล่งเสียงเคลื่อนที่เข้าหาผู้สังเกตที่อยู่กับที่
\item ค) เมื่อทั้งสองเคลื่อนที่ห่างออกจากกัน
\item ง) เมื่อทั้งสองอยู่กับที่
\end{enumerate}
\item \textbf{คลื่นเสียงไม่สามารถเดินทางผ่านข้อใดต่อไปนี้ได้}
\begin{enumerate}
\item ก) น้ำ
\item ข) อากาศ
\item ค) เหล็ก
\item ง) สุญญากาศ
\end{enumerate}
\item \textbf{ลำโพงสองตัวมีกำลัง 10 W และ 40 W ที่ระยะห่างเท่ากัน ระดับเสียงต่างกันกี่ dB}
\begin{enumerate}
\item ก) 2 dB
\item ข) 4 dB
\item ค) 6 dB
\item ง) 10 dB
\end{enumerate}
\item \textbf{เรือดำน้ำส่งคลื่นเสียงความถี่ 1500 Hz ลงไปในน้ำทะเล คลื่นสะท้อนกลับจากปลาวาฬแล้วรับได้ที่ความถี่ 1580 Hz ปลาวาฬเคลื่อนที่เข้าหาเรือดำน้ำหรือไม่ ด้วยความเร็วเท่าใด (v = 1500 m/s)}
\begin{enumerate}
\item ก) เข้าหา, ประมาณ 26 m/s
\item ข) เข้าหา, ประมาณ 78 m/s
\item ค) ห่างออก, ประมาณ 26 m/s
\item ง) ห่างออก, ประมาณ 78 m/s
\end{enumerate}
\item \textbf{ช่วงความถี่ใดที่เสียงพูดของมนุษย์อยู่}
\begin{enumerate}
\item ก) 100 Hz - 500 Hz
\item ข) 300 Hz - 3,000 Hz
\item ค) 500 Hz - 5,000 Hz
\item ง) 2,000 Hz - 10,000 Hz
\end{enumerate}
\item \textbf{ระดับเสียง 70 dB เพิ่มเป็น 90 dB ความเข้มเพิ่มขึ้นกี่เท่า}
\begin{enumerate}
\item ก) 2 เท่า
\item ข) 20 เท่า
\item ค) 100 เท่า
\item ง) 1000 เท่า
\end{enumerate}
\item \textbf{แหล่งเสียงความถี่ 800 Hz เคลื่อนที่ห่างออกจากผู้สังเกตด้วยความเร็ว 20 m/s ความถี่ที่ได้ยินเป็นเท่าใด (v = 340 m/s)}
\begin{enumerate}
\item ก) 700 Hz
\item ข) 750 Hz
\item ค) 760 Hz
\item ง) 800 Hz
\end{enumerate}
\item \textbf{อัตราเร็วเสียงในอากาศจะเปลี่ยนแปลงอย่างไรเมื่ออุณหภูมิเพิ่มขึ้น}
\begin{enumerate}
\item ก) ลดลง
\item ข) คงที่
\item ค) เพิ่มขึ้น
\item ง) เพิ่มแล้วลด
\end{enumerate}
\item \textbf{เครื่องเล่นเพลงระดับเสียง 85 dB กับเครื่องอีกเครื่องที่ 95 dB ความเข้มต่างกันกี่เท่า}
\begin{enumerate}
\item ก) 2 เท่า
\item ข) 5 เท่า
\item ค) 10 เท่า
\item ง) 100 เท่า
\end{enumerate}
\item \textbf{รถพยุมยางเสียงดัง 440 Hz วิ่งเข้าหาทางโค้งด้วยความเร็ว 30 m/s ระยะห่างจากทางโค้ง 200 m ความถี่ที่ได้ยินเป็นเท่าใด (v = 340 m/s)}
\begin{enumerate}
\item ก) 440 Hz
\item ข) 480 Hz
\item ค) 500 Hz
\item ง) 520 Hz
\end{enumerate}
\item \textbf{ถ้าอัตราเร็วเสียงในอากาศ = 340 m/s และความถี่ = 500 Hz ความยาวคลื่นเป็นเท่าใด}
\begin{enumerate}
\item ก) 0.34 m
\item ข) 0.68 m
\item ค) 1.36 m
\item ง) 170 m
\end{enumerate}
\item \textbf{เครื่องบินบินด้วยความเร็ว 600 m/s ซึ่งมากกว่าความเร็วเสียง (Mach 1.76) ความถี่ที่คนบนพื้นได้ยินจะเป็นอย่างไร}
\begin{enumerate}
\item ก) ได้ยินปกติ
\item ข) ได้ยินเสียงโซนิคบูม
\item ค) ไม่ได้ยินเลย
\item ง) ได้ยินเสียงต่ำกว่าปกติ
\end{enumerate}
\end{enumerate}
\end{document}